% ---
% titre: Problème de proportionnalité - la vente de pâtisserie
% theme: FA
% chapitre: Proportionnalité
% sous_chapitre: Résoudre des problèmes
% niveau: [10e]
% type: EVAL
% tags: [proportionnalite,c-nombres-decimaux, p-question-TS]
% difficulte: 1
% corrige: true
% ---

\question[2]
Au marché, 5 barquettes de fraises coûtent 5,75 francs.

\begin{parts}
\vspace{20pt}\part[2] Calcule le prix de 4 barquettes de fraises.

\vspace{10pt}{\footnotesize\raggedleft Espace pour ta démarche et tes calculs\par}
\begin{solutionorgrid}[140pt]

\vspace{5pt}
\def\arraystretch{1.5}
\begin{tabularx}{0.5\textwidth}{l|C|C|}
Prix [CHF]&5,75&\\
\hline
Nombre de barquettes&5&4\\
\end{tabularx}
\def\arraystretch{1.5}

\hfill\textbf{(0.5pt pour la construction du tableau)}

\vspace{15pt}$4 \cdot 5,75 : 5 = 4,60 francs$ \hfill calcul \textbf{(0,5pt) et réponse \textbf{(0,5pt)}

\end{solutionorgrid}

\begin{tcolorbox}[enhanced jigsaw, interior hidden, colframe=box, parbox=false]%
Réponse: 
\sol{\vspace{20pt}}{Quatre barquettes coûtent 4,60 francs. \textbf{(0.5pt)}}
\end{tcolorbox}

\vspace{20pt}\part[2] Combien peut-on acheter de barquettes avec 12,65 francs?

\vspace{10pt}{\footnotesize\raggedleft Espace pour ta démarche et tes calculs\par}
\begin{solutionorgrid}[140pt]

\vspace{5pt}
\def\arraystretch{1.5}
\begin{tabularx}{0.5\textwidth}{l|C|C|}
Prix [CHF]&5,75&12,65\\
\hline
Nombre de barquettes&5&\\
\end{tabularx}
\def\arraystretch{1.5}

\hfill\textbf{(0.5pt pour la construction du tableau)}

\vspace{15pt}$5 \cdot 12,65 : 5,75 = 11 barquettes$ \hfill calcul \textbf{(0,5pt) et réponse \textbf{(0,5pt)}

\end{solutionorgrid}

\begin{tcolorbox}[enhanced jigsaw, interior hidden, colframe=box, parbox=false]%
Réponse: 
\sol{\vspace{20pt}}{Je peux acheter 11 barquettes avec 12,65. \textbf{(0.5pt)}}
\end{tcolorbox}



\end{parts}


